\documentclass[12pt]{article}
\usepackage{fontspec}
\usepackage{xunicode}
\usepackage{polyglossia}
\setmainlanguage{french}
%Avec polyglossia et le français, une espace fine est ajoutée automatiquement devant « ! » et « ; ».
%Mettez les deux lignes ci-dessus en commentaire et recompilez pour voir la différence.

\begin{document}

	Ceci est un exemple de ligne qui se termine avec une césure automatiquement placée.
	Insérez      des     espaces    où vous    voulez
	entre
	les mots.
	%Plusieurs espaces à la suite, ou un simple retour à la ligne, comptent pour une seule espace : le rendu ne change pas.

	%Une ligne blanche (ci-dessus) crée un nouveau paragraphe.
	Essayez de faire un paragraphe au début de cette phrase. Mettez un espace~insécable après le mot "espace", et observez la ponctuation: ici! et là; ou là.
	%Le tilde ~ est une espace insécable : « espace » et « insécable » ne seront jamais séparés par un saut de ligne.

	Voici différents tirets: tiret simple (-) pour les mots composés; demi-cadratin (--) pour les plages de nombre; cadratin (---) pour les incises.

	Un sous-titre s'écrit avec un tiret simple, les pages 12--15 avec un demi-cadratin.



	Utilisez des cadratins à la place des parenthèses --- c'est plus élégant --- dans votre mémoire.
	%Plusieurs lignes blanches à la suite ne créent qu'un seul saut de paragraphe.

	%Cet exercice sert à montrer comment taper du texte
	%Ceci est un commentaire: apparaît-il?
	%Non : tout ce qui suit un % est ignoré à la compilation.

\end{document}
